\documentclass[12pt]{article}
\usepackage[utf8]{inputenc}
\usepackage[T1]{fontenc}
\usepackage{hyperref}
\title{Wemale Language Resources for K-12 Curriculum Development}
\author{Compiled by Claude Sonnet 5 (claude-sonnet-5) for Andr\'as Kornai}
\date{}
\begin{document}
\maketitle
\section{Introduction}
Wemale (ISO 639-3 weo) is an Austronesian language of western Seram, Maluku, Indonesia, with northern and southern varieties. Speaker estimates vary widely, from about 5,000 to 11,000 depending on the source. Most speakers also use Indonesian for school and trade. Searches found no educational materials.
\section{Primary Educational Resources}
None found.
\section{Linguistic and Reference Materials}
\begin{itemize}
\item Wemale language, Wikipedia: \url{https://en.wikipedia.org/wiki/Wemale_language}
\item Wemale people, Wikipedia: \url{https://en.wikipedia.org/wiki/Wemale_people}
\end{itemize}
\section{Religious and Community Resources}
\begin{itemize}
\item Joshua Project profile, North Wemale: \url{https://joshuaproject.net/people_groups/15907}
\item Bible-society language record ``697 weo 001'' (title as listed; the content was not checked): \url{https://bakat.sabda.org/697-weo-001.html}
\end{itemize}
\section{Digital and Audio Resources}
None found.
\section{Curriculum Development Context}
The language is endangered and used by adults only, according to the brief. No orthography, primer or school programme was found. Any curriculum would start from documentation and community work, with Indonesian as the school language of support.
\end{document}